The PD of a grade under informative rating withdrawals is a curve in the odds ratio $k$ of withdrawal, not a number, and the three treatments in use are three points of it. The curve, its worst-case bound and the tipping $k$ of a decision come in closed form from the transition page alone; a second source that counts the defaults of the same ratings, such as the default page of the ESMA repository, anchors the curve and gives an interval. On rating paths of known truth the anchored interval covers the true PD of grade B in \nv{h1.min}\% to 100\% of \nv{reps} repetitions at informativeness 1, 3 and 9 in the first-year and pooled designs and at 9 in the ten-fold design, against \nv{cov.pool.i9.industry}\% for removal in the pooled design, and it is \nv{w.fy.i9.anchored} percentage points wide in the first annual cohort against \nv{w.fy.i9.followed} for the oracle that follows every withdrawal, which it contains by construction. Its centre is exact when the second source counts exactly what the transition page would count and carries the capture of the second source as a bias otherwise. On the agencies' cohorts, the pooled anchored $k$ is above 1 in every grade, \nv{real.sp.CCC-C.k} for S\&P's CCC to C, but it does not hold across years for S\&P's grades B and CCC to C or for Fitch's grade B, and the instability has the size of a definition gap that changes from year to year, not of sampling. The anchored odds ratio of the agencies is conditional on the two pages defining default alike. The next steps are an interval that treats the noise of the two sources jointly, to narrow the conservative width, a bank's own second source in place of the repository's pages, and a multi-year extension that models the selection of pooled cohorts.
